\documentclass[12pt,a4paper]{report}

% ── Codificação e Idioma ──────────────────────────────────────────────────────
\usepackage[utf8]{inputenc}
\usepackage[T1]{fontenc}
\usepackage[brazil]{babel}

% ── Tipografia ────────────────────────────────────────────────────────────────
\usepackage{lmodern}
\usepackage{microtype}

% ── Layout ───────────────────────────────────────────────────────────────────
\usepackage[top=3cm,bottom=2cm,left=3cm,right=2cm]{geometry}
\usepackage{setspace}
\onehalfspacing

% ── Gráficos e Figuras ────────────────────────────────────────────────────────
\usepackage{graphicx}
\usepackage{float}
\usepackage{tikz}
\usetikzlibrary{arrows.meta,shapes.geometric,positioning,fit,backgrounds}

% ── Tabelas ───────────────────────────────────────────────────────────────────
\usepackage{booktabs}
\usepackage{tabularx}
\usepackage{longtable}
\usepackage{multirow}

% ── Código ────────────────────────────────────────────────────────────────────
\usepackage{listings}
\usepackage{xcolor}

\lstset{
  basicstyle=\ttfamily\small,
  keywordstyle=\color{blue}\bfseries,
  commentstyle=\color{gray}\itshape,
  stringstyle=\color{teal},
  breaklines=true,
  frame=single,
  numbers=left,
  numberstyle=\tiny\color{gray},
  captionpos=b,
  showstringspaces=false,
}

% ── Matemática ───────────────────────────────────────────────────────────────
\usepackage{amsmath,amssymb,mathtools}
\DeclareMathOperator*{\argmax}{arg\,max}
\DeclareMathOperator*{\argmin}{arg\,min}

% ── Hiperlinks ───────────────────────────────────────────────────────────────
\usepackage[hidelinks]{hyperref}

% ── Referências ──────────────────────────────────────────────────────────────
\usepackage[style=abnt,backend=biber]{biblatex}
\addbibresource{referencias.bib}

% ── Cabeçalho e Rodapé ───────────────────────────────────────────────────────
\usepackage{fancyhdr}
\pagestyle{fancy}
\fancyhf{}
\fancyhead[R]{\small\textit{FasiTech — Fundamentos Científicos}}
\fancyfoot[C]{\thepage}

% ── Caixas de Destaque ────────────────────────────────────────────────────────
\usepackage{tcolorbox}
\tcbuselibrary{skins,breakable}

\newtcolorbox{insight}[1][]{
  colback=blue!5!white,colframe=blue!40!black,
  fonttitle=\bfseries, title=Decisão de Projeto,#1
}

% ─────────────────────────────────────────────────────────────────────────────
\begin{document}

\begin{titlepage}
  \centering
  \vspace*{2cm}
  {\Large\bfseries Universidade Federal do Pará\\
  Faculdade de Sistemas de Informação\par}
  \vspace{3cm}
  {\Huge\bfseries FasiTech\\[0.5em]
  \large Fundamentos Científicos do Sistema\par}
  \vspace{2cm}
  {\large
  Módulo de Recuperação com Geração Aumentada (RAG)\\
  Módulo de Extração Inteligente de Documentos (ACC)
  \par}
  \vfill
  {\large Belém -- PA\\2026}
\end{titlepage}

\tableofcontents
\clearpage

% ─────────────────────────────────────────────────────────────────────────────
\chapter{Introdução}
% ─────────────────────────────────────────────────────────────────────────────

O sistema FasiTech é uma plataforma web institucional desenvolvida para centralizar
fluxos acadêmicos da Faculdade de Sistemas de Informação (FASI) da UFPA.
Entre seus módulos, dois se apoiam diretamente em técnicas de Inteligência Artificial:

\begin{enumerate}
  \item \textbf{Diretor Virtual} — assistente conversacional baseado em
        \textit{Retrieval-Augmented Generation} (RAG) que responde perguntas
        sobre regulamentos, cronogramas e documentos acadêmicos.
  \item \textbf{Módulo ACC} — extrator automático de cargas horárias de
        certificados enviados em formato PDF, usando Modelos de Linguagem
        Multimodais (\textit{Multimodal Large Language Models}).
\end{enumerate}

Este documento registra a ciência aplicada em cada componente, servindo de
referência para manutenção, evolução e publicação futura dos resultados.

% ─────────────────────────────────────────────────────────────────────────────
\chapter{Módulo RAG — Diretor Virtual}
% ─────────────────────────────────────────────────────────────────────────────

\section{Fundamentos: Retrieval-Augmented Generation}

\textit{Retrieval-Augmented Generation} (RAG) é uma arquitetura que combina
recuperação de documentos com geração de linguagem natural.
A formulação original foi proposta por \textcite{lewis2020rag}:
dado um modelo generativo \(p_\theta(y \mid x)\) e um índice de documentos $\mathcal{D}$,
o modelo recupera os documentos mais relevantes \(\mathbf{z} = \text{Top-}k(q, \mathcal{D})\)
e condiciona a geração sobre eles:

\begin{equation}
  p(y \mid x) = \sum_{z \in \mathbf{z}} p_\eta(z \mid x)\, p_\theta(y \mid x, z)
  \label{eq:rag}
\end{equation}

onde $p_\eta(z \mid x)$ é a probabilidade do documento $z$ ser relevante para
a consulta $x$, calculada via busca densa, e $p_\theta$ é o LLM gerador.

\begin{insight}[title=Por que RAG e não fine-tuning?]
Fine-tuning parametriza o conhecimento no modelo, tornando atualizações
custosas e sem rastreabilidade de fonte. RAG externaliza o conhecimento
nos documentos, permitindo atualização imediata sem retreinamento e
garantindo que respostas sejam fundamentadas em trechos verificáveis.
\end{insight}

\subsection{Recuperação Densa com Bi-Encoders}

A recuperação densa (DR — \textit{Dense Retrieval}) representa consultas e
documentos como vetores densos em um espaço semântico compartilhado.
O trabalho seminal de \textcite{karpukhin2020dpr} (DPR —
\textit{Dense Passage Retrieval}) demonstrou que bi-encoders treinados por
aprendizado contrastivo superam a busca esparsa BM25 em domínios factuais.

A similaridade entre a consulta $q$ e um documento $d$ é calculada por:

\begin{equation}
  \text{sim}(q, d) = \cos\!\left(E_Q(q),\, E_D(d)\right)
  = \frac{E_Q(q) \cdot E_D(d)}{\|E_Q(q)\|\,\|E_D(d)\|}
  \label{eq:cosine}
\end{equation}

No sistema FasiTech, o embedder é o \textbf{Ollama Embedder}
(\texttt{nomic-embed-text} \autocite{nussbaum2024nomic}), executado localmente via servidor Ollama
\autocite{ollama2023}
(\texttt{OLLAMA\_HOST}, padrão \texttt{localhost:11434}), produzindo
vetores de 768 dimensões. Todos os vetores são normalizados pelo
serviço (norma L2 = 1) antes de serem armazenados no LanceDB,
permitindo que a similaridade cosseno seja computada como produto
escalar simples — $\text{sim}(q,d) = E(q) \cdot E(d)$.

\subsection{Busca Híbrida: Semântica + BM25}

Busca híbrida combina recuperação densa (semântica) com recuperação esparsa
baseada em frequência de termos \autocite{robertson1994okapi}.
O score BM25 para um documento $d$ dado uma consulta $q$ com termos $\{q_i\}$ é:

\begin{equation}
  \text{BM25}(d,q) = \sum_{i} \text{IDF}(q_i) \cdot
  \frac{f(q_i,d)\,(k_1+1)}{f(q_i,d) + k_1\!\left(1 - b + b\,\frac{|d|}{\text{avgdl}}\right)}
  \label{eq:bm25}
\end{equation}

onde $k_1$ e $b$ são hiperparâmetros de saturação de frequência e normalização
por comprimento do documento, respectivamente.

O banco vetorial \textbf{LanceDB} utilizado no FasiTech configura
\texttt{SearchType.hybrid}, combinando os dois rankings via
\textit{Reciprocal Rank Fusion} (RRF) \autocite{cormack2009rrf}:

\begin{equation}
  \text{RRF}(d) = \sum_{r \in R} \frac{1}{k + r(d)}
  \label{eq:rrf}
\end{equation}

onde $r(d)$ é a posição do documento no ranking $r$ e $k=60$ é constante.
Além disso, o módulo \texttt{\_keyword\_search} realiza busca adicional por
palavras-chave como complemento pós-semântico.

\section{Estratégias de Chunking}

\subsection{Chunking Básico com Sobreposição}

A fragmentação de documentos (\textit{chunking}) é etapa crítica no pipeline
RAG. O método original no FasiTech divide textos em blocos de 800 caracteres
com sobreposição (\textit{overlap}) de 150 caracteres, respeitando quebras
de parágrafo. Esta abordagem, denominada \textit{fixed-size chunking with
overlap}, minimiza a perda de contexto nas fronteiras dos chunks.

\subsection{Chunking Contextual de Seção (Melhoria 1)}
\label{sec:melhoria1}

A limitação do chunking simples é a perda do contexto hierárquico do documento.
Quando um chunk pertence à seção ``Período Letivo 2026.2'', o título dessa
seção pode estar em um chunk anterior, tornando a recuperação dependente de
que o usuário mencione explicitamente ``2026.2'' na consulta.

Inspirado no \textit{Contextual Retrieval} proposto pela Anthropic
\autocite{anthropic2024contextual} e no trabalho de \textcite{shi2023large}
sobre prepend de contexto para melhoria de recuperação, o FasiTech implementa
o método \texttt{\_split\_with\_section\_context}:

\begin{enumerate}
  \item O documento Markdown é percorrido linha a linha.
  \item Cabeçalhos (\texttt{\#}, \texttt{\#\#}, \texttt{\#\#\#},
        \texttt{\#\#\#\#}) são detectados via expressão regular
        \texttt{\^{}(\#\{1,4\})\textbackslash{}s+(.+)} e armazenados
        como \texttt{section\_header} corrente.
  \item Cada chunk gerado pelo método base (\texttt{\_split\_into\_chunks})
        recebe um \emph{prefixo contextual}:
        \begin{center}
          \texttt{[Documento: X | Seção: Y]$\backslash$n$\backslash$n<chunk>}
        \end{center}
  \item O par \texttt{(section\_header, chunk\_contextualizado)} é retornado;
        o campo \texttt{section} é persistido nos metadados do LanceDB.
  \item O chunk contextualizado — não o texto bruto — é embeddado e indexado.
\end{enumerate}

\begin{lstlisting}[language=Python,
  caption={Núcleo de \texttt{\_split\_with\_section\_context} (\texttt{rag\_ppc.py}).}]
lines = text.split("\n")
heading_re = re.compile(r"^(#{1,4})\s+(.+)")

# Agrupa linhas por secao
sections: List[tuple[str, List[str]]] = []
current_header, current_lines = "", []
for line in lines:
    m = heading_re.match(line)
    if m:
        if current_lines:
            sections.append((current_header, current_lines))
        current_header = m.group(2).strip()
        current_lines = []
    else:
        current_lines.append(line)

# Gera chunks contextualizados
result: List[tuple[str, str]] = []
for sec_header, sec_lines in sections:
    block = "\n".join(sec_lines).strip()
    raw_chunks = ChatbotService._split_into_chunks(block, chunk_size, overlap)
    for chunk in raw_chunks:
        prefix = f"[Documento: {doc_name} | Secao: {sec_header}]\n\n"
        contextualized = (prefix + chunk) if sec_header else chunk
        result.append((sec_header, contextualized))
\end{lstlisting}

O efeito é que o embedding do chunk carrega informação da seção pai,
aumentando a similaridade com consultas que referenciam aquela seção,
mesmo que o chunk em si não a mencione.

\begin{figure}[H]
\centering
\begin{tikzpicture}[
  node distance=0.6cm and 1.2cm,
  box/.style={draw,rounded corners,fill=blue!8,minimum width=3.8cm,minimum height=0.8cm,align=center,font=\small},
  arrow/.style={-{Stealth},thick}
]
  \node[box] (doc)  {Documento Markdown};
  \node[box,below=of doc] (parse) {Parse de Headings\\(\texttt{\#\#} etc.)};
  \node[box,below=of parse] (chunk) {Chunking por\\Parágrafo};
  \node[box,below=of chunk] (prefix) {Prefixo:\\{[Doc | Seção]}};
  \node[box,below=of prefix] (embed) {Embedding\\(Gemini 768-d)};
  \node[box,below=of embed] (lance) {LanceDB\\(Hybrid Index)};
  \draw[arrow] (doc) -- (parse);
  \draw[arrow] (parse) -- (chunk);
  \draw[arrow] (chunk) -- (prefix);
  \draw[arrow] (prefix) -- (embed);
  \draw[arrow] (embed) -- (lance);
\end{tikzpicture}
\caption{Pipeline de indexação com chunking contextual de seção.}
\label{fig:chunking}
\end{figure}

\subsection{Chunking Pai-Filho (\textit{Small-to-Big})}

Como direção futura, o FasiTech pode adotar o \textit{Parent-Child Chunking}
introduzido no LlamaIndex \autocite{llamaindex2023}: chunks pequenos
(200 chars) são indexados para busca precisa, mas o contexto entregue ao
LLM é o chunk pai (1\,200 chars) que contém o trecho recuperado.
Isso reduz o ruído na busca vetorial enquanto mantém contexto suficiente
para a geração.

\section{Cache Semântico}

\subsection{Motivação e Arquitetura}

O cache semântico armazena pares (pergunta, resposta) avaliados positivamente
e os serve em consultas futuras semanticamente similares, eliminando chamadas
ao LLM para perguntas recorrentes. Esta abordagem é referenciada em
\textcite{bang2023gptcache} e \textcite{zhu2023semsimilaritycache}.

O sistema FasiTech implementa um cache em duas camadas no LanceDB:

\begin{table}[H]
\centering
\caption{Parâmetros do cache semântico do FasiTech.}
\label{tab:cache}
\begin{tabularx}{\textwidth}{lXl}
\toprule
\textbf{Parâmetro} & \textbf{Descrição} & \textbf{Valor} \\
\midrule
Threshold de similaridade & Cosseno mínimo para hit de cache & 0,90 \\
Avaliação mínima (\textit{trusted}) & Média de notas para promover entrada & 4,4/5 \\
Avaliações para promoção & Qtd. mínima de feedbacks & 2 \\
TTL \textit{trusted} & Validade de entrada confiável & 30 dias \\
TTL \textit{candidate} & Validade de entrada candidata & 14 dias \\
\bottomrule
\end{tabularx}
\end{table}

\subsection{Ciclo de Vida das Entradas}

Uma entrada percorre o seguinte ciclo:

\begin{enumerate}
  \item \textbf{Candidate} — criada após primeira avaliação positiva (nota $\geq 3$).
  \item \textbf{Trusted} — promovida quando $\bar{r} \geq 4{,}4$ com $n \geq 2$ avaliações.
  \item \textbf{Servida} — retornada diretamente ao usuário sem chamar o LLM.
  \item \textbf{Expirada} — removida após TTL ou se nota $\leq 2$ na última avaliação.
\end{enumerate}

A função de confiança $C$ combina qualidade e volume:

\begin{equation}
  C = \frac{\bar{r} - 1}{4} \cdot \min\!\left(1,\,\frac{n}{n_{\min}}\right)
\end{equation}

Adicionalmente, cada entrada armazena o \texttt{documents\_hash}, garantindo
que mudanças nos documentos fontes invalidem automaticamente o cache.

\section{Filtragem de Domínio}

Antes de acionar o RAG, o sistema verifica se a pergunta pertence ao domínio
acadêmico para evitar respostas sobre tópicos fora de escopo.

A estratégia é bifásica: (1) detecção de palavras-chave acadêmicas
(ex.: ``TCC'', ``matrícula'', ``ACC'') ajusta o limiar de similaridade para
mais permissivo ($\tau = 0{,}50$); (2) na ausência de palavras-chave, um
limiar mais estrito é aplicado ($\tau = 0{,}72$) para evitar que perguntas
genéricas em português sejam processadas por similaridade superficial de idioma.

O score combinado usa 40\% da melhor similaridade e 60\% da média dos
três primeiros resultados:

\begin{equation}
  S = 0{,}4 \cdot \text{sim}_{\max} + 0{,}6 \cdot \bar{\text{sim}}_3
\end{equation}

\section{Detecção de Ambiguidade Temporal (Melhoria 2)}
\label{sec:melhoria2}

Estudos em busca conversacional demonstram que perguntas ambíguas — aquelas
que carecem de entidades contextuais necessárias para uma resposta precisa —
devem ser respondidas com uma pergunta de esclarecimento antes da recuperação
\autocite{aliannejadi2021clariq}.

No contexto do FasiTech, a ambiguidade temporal ocorre quando o usuário
pergunta sobre prazos sem especificar o período letivo. O método
\texttt{\_needs\_temporal\_clarification} implementa a seguinte lógica:

\begin{equation}
  \text{clarify}(q) = \underbrace{\exists\, t \in T_{\text{trigger}} : t \in q}_{\text{gatilho temporal}}
  \;\wedge\;
  \underbrace{\nexists\, e \in E_{\text{period}} : e \in q}_{\text{ausência de entidade}}
\end{equation}

Os conjuntos são implementados como \texttt{frozenset} para busca $O(1)$:

\begin{lstlisting}[language=Python,
  caption={Constantes de detecção temporal (\texttt{rag\_ppc.py}).}]
_TEMPORAL_TRIGGERS: frozenset = frozenset({
    "periodo", "prazo", "submissao", "entrega", "data", "quando",
    "limite", "cronograma", "calendario", "vencimento", "ate quando",
    "deadline", "datas", "prazos", "agenda",
})
_TEMPORAL_ENTITIES: frozenset = frozenset({
    # Identificadores de periodo letivo: 2024.1 ... 2028.4
    "2024.1", "2024.2", "2024.3", "2024.4",
    "2025.1", "2025.2", "2025.3", "2025.4",
    "2026.1", "2026.2", "2026.3", "2026.4",
    "2027.1", ..., "2028.4",
    # Expressoes naturais de periodo
    "primeiro periodo", "segundo periodo", "semestre atual",
    "proximo semestre", "este semestre",
})

@staticmethod
def _needs_temporal_clarification(question: str) -> bool:
    q_lower = question.lower()
    has_trigger = any(t in q_lower for t in _TEMPORAL_TRIGGERS)
    has_entity  = any(e in q_lower for e in _TEMPORAL_ENTITIES)
    return has_trigger and not has_entity
\end{lstlisting}

Quando $\text{clarify}(q) = \top$, o método \texttt{\_get\_available\_periods}
varre todos os chunks do LanceDB em busca do padrão
\texttt{\textbackslash b(20\textbackslash d\{2\}.\textbackslash{}[1-4\textbackslash{}])\textbackslash b}
e retorna os períodos reais indexados:

\begin{lstlisting}[language=Python,
  caption={Descoberta dinâmica de períodos (\texttt{rag\_ppc.py}).}]
def _get_available_periods(self) -> List[str]:
    period_re = re.compile(r"\b(20\d{2}\.[1-4])\b")
    found: set = set()
    df = lancedb.connect(self.db_url).open_table("recipes").to_pandas()
    for row in df.itertuples():
        p = json.loads(row.payload)
        for m in period_re.findall(p.get("content", "")):
            found.add(m)
    return sorted(found)   # ex.: ['2026.2', '2026.3', '2026.4']
\end{lstlisting}

A resposta de esclarecimento é formatada em Markdown (renderizado no
frontend via \texttt{react-markdown} + \texttt{remark-gfm}):

\begin{lstlisting}[language=Python,
  caption={Mensagem de esclarecimento gerada dinamicamente.}]
period_list = ", ".join(available_periods)  # ex.: 2026.2, 2026.3, 2026.4
clarification = (
    f"Para qual **periodo letivo** voce deseja as informacoes?\n\n"
    f"Periodos disponiveis nos documentos: **{period_list}**\n\n"
    f"Por favor, especifique o periodo (ex.: *2026.2*) para que eu "
    f"possa te dar uma resposta precisa."
)
\end{lstlisting}

Esta abordagem elimina chamadas desnecessárias ao LLM antes de
executar o pipeline RAG completo com contexto insuficiente.

\begin{insight}[title=Ganho esperado]
Esta abordagem elimina a necessidade de o usuário reformular a pergunta.
Em vez de receber uma resposta incompleta ou mesclada entre períodos,
o sistema pede esclarecimento na primeira interação, reduzindo o número
médio de turnos para resposta de $\approx 2$ para $1$.
\end{insight}

\section{Expansão de Consulta por Siglas}

O método \texttt{\_expand\_query} implementa um mapa de siglas para seus
expandidos semânticos (ex.: ``ACC'' $\to$ ``Atividade Curricular
Complementar atividades complementares ACG extracurricular'').

Esta técnica é variante do \textit{Query Expansion} estudado em
\textcite{carpineto2012survey}: ao adicionar termos sinônimos à consulta
antes do embedding, aumenta-se o recall da busca vetorial sem modificar
o banco de vetores.

\section{Visão Geral do Pipeline RAG}

\begin{figure}[H]
\centering
\begin{tikzpicture}[
  node distance=0.5cm and 0.4cm,
  proc/.style={draw,rounded corners,fill=green!8,minimum width=4.0cm,
               minimum height=0.75cm,align=center,font=\footnotesize},
  dec/.style={draw,diamond,fill=orange!12,aspect=2,
              minimum width=3.2cm,minimum height=0.6cm,
              align=center,font=\footnotesize},
  term/.style={draw,rounded corners,fill=gray!12,minimum width=3.2cm,
               minimum height=0.6cm,align=center,font=\footnotesize},
  arrow/.style={-{Stealth},thick}
]
  \node[proc]  (q)     {Pergunta do Usuário};
  \node[dec,below=of q]  (cache)  {Cache\\Hit?};
  \node[term,right=2.8cm of cache] (ret)    {Retorna do Cache};
  \node[dec,below=0.8cm of cache]  (temp)   {Ambiguidade\\Temporal?};
  \node[term,right=2.8cm of temp]  (clar)   {Pergunta de\\Esclarecimento};
  \node[dec,below=0.8cm of temp]   (dom)    {In-Domain?};
  \node[term,right=2.8cm of dom]   (ood)    {Fora de Escopo};
  \node[proc,below=0.8cm of dom]   (exp)    {Expansão de Siglas};
  \node[proc,below=of exp]         (retr)   {Busca Híbrida\\(Semântica + BM25)};
  \node[proc,below=of retr]        (prompt) {Injeção de Contexto\\no Prompt};
  \node[proc,below=of prompt]      (llm)    {LLM (Gemini 2.5 Flash\\ou Ollama fallback)};
  \node[proc,below=of llm]         (ans)    {Resposta ao Usuário};

  \draw[arrow] (q) -- (cache);
  \draw[arrow] (cache) -- node[right,font=\tiny]{Não} (temp);
  \draw[arrow] (cache) -- node[above,font=\tiny]{Sim} (ret);
  \draw[arrow] (temp) -- node[right,font=\tiny]{Não} (dom);
  \draw[arrow] (temp) -- node[above,font=\tiny]{Sim} (clar);
  \draw[arrow] (dom) -- node[right,font=\tiny]{Sim} (exp);
  \draw[arrow] (dom) -- node[above,font=\tiny]{Não} (ood);
  \draw[arrow] (exp) -- (retr);
  \draw[arrow] (retr) -- (prompt);
  \draw[arrow] (prompt) -- (llm);
  \draw[arrow] (llm) -- (ans);
\end{tikzpicture}
\caption{Fluxo completo do pipeline RAG do Diretor Virtual.}
\label{fig:pipeline}
\end{figure}

% ─────────────────────────────────────────────────────────────────────────────
\chapter{Módulo ACC — Extração Inteligente de Documentos}
% ─────────────────────────────────────────────────────────────────────────────

\section{Contexto e Desafio}

As Atividades Complementares de Graduação (ACC/ACG) são requisito curricular
para integralização do curso de Sistemas de Informação.
Historicamente, a validação dessas atividades é feita manualmente por
servidores, que examinam cada certificado entregue pelo aluno e somam
as cargas horárias.

O processo é:
\begin{itemize}
  \item \textbf{Heterogêneo} — certificados variam em formato, layout, idioma e fonte.
  \item \textbf{Propenso a erros} — leitura manual de dezenas de certificados por aluno.
  \item \textbf{Não escalável} — custo cresce linearmente com o número de alunos.
\end{itemize}

O módulo ACC do FasiTech automatiza a extração da carga horária via
\textit{Document AI} com Modelos de Linguagem Multimodais.

\section{Fundamentos: Document AI e Multimodal LLMs}

\textit{Document AI} é o campo de estudo da compreensão automática de
documentos não-estruturados \autocite{cui2021documentai}.
Diferentemente de NLP clássico — que opera sobre texto puro —,
Document AI lida com a informação visual e estrutural de documentos
(layout, tabelas, fontes, posição do texto).

\subsection{Modelos Multimodais}

Modelos de linguagem multimodais (\textit{Multimodal LLMs}) processam
imagens e texto de forma conjunta, sem necessidade de OCR separado.
O \textbf{Gemini} \autocite{team2023gemini} é um modelo multimodal
nativo que interpreta texto, imagens e PDFs em uma única passagem,
sendo especialmente adequado para certificados escaneados.

A geração de resposta $y$ dado um documento $D$ (página PDF) e uma
instrução $I$ pode ser expressa como:

\begin{equation}
  y = \argmax_y\; p_\theta(y \mid D, I)
  \label{eq:multimodal}
\end{equation}

onde $D$ é representado como sequência de tokens visuais extraídos por
um \textit{Vision Encoder} e $I$ é o prompt de instrução.

\subsection{Processamento de PDFs com pdf2image e Poppler}

Quando o modelo não suporta PDF nativo, a conversão para imagens de alta
resolução (300 DPI) via \textit{pdf2image}/\textit{Poppler} é prática
padrão \autocite{blecher2023nougat}. O FasiTech oferece suporte a ambos
os modos:

\begin{enumerate}
  \item \textbf{PDF direto} — enviado como objeto \texttt{File} ao Gemini,
        que processa nativamente todas as páginas.
  \item \textbf{PDF $\to$ imagem} — cada página é convertida a PNG 300 DPI
        e enviada individualmente ao modelo via \texttt{Image}.
\end{enumerate}

\section{Arquitetura do Módulo ACC}

\begin{figure}[H]
\centering
\begin{tikzpicture}[
  node distance=0.6cm,
  proc/.style={draw,rounded corners,fill=teal!8,minimum width=4.2cm,
               minimum height=0.8cm,align=center,font=\small},
  arrow/.style={-{Stealth},thick}
]
  \node[proc] (upload)  {Upload do PDF\\(certificados mesclados)};
  \node[proc,below=of upload] (val)     {Validação\\(tamanho, formato)};
  \node[proc,below=of val]    (agent)   {Agente Gemini\\(GEMINI\_MODEL\_VISION)};
  \node[proc,below=of agent]  (prompt)  {Prompt Estruturado\\(extração por página)};
  \node[proc,below=of prompt] (parse)   {Parse da Resposta\\(regex TOTAL GERAL)};
  \node[proc,below=of parse]  (save)    {Persistência\\(BD + Google Sheets)};
  \node[proc,below=of save]   (email)   {Notificação\\por E-mail};

  \draw[arrow] (upload) -- (val);
  \draw[arrow] (val) -- (agent);
  \draw[arrow] (agent) -- (prompt);
  \draw[arrow] (prompt) -- (parse);
  \draw[arrow] (parse) -- (save);
  \draw[arrow] (save) -- (email);
\end{tikzpicture}
\caption{Pipeline de processamento do Módulo ACC.}
\label{fig:acc}
\end{figure}

\section{Prompt Engineering para Extração Estruturada}

O desempenho de LLMs em tarefas de extração de informação depende
criticamente da qualidade do prompt \autocite{brown2020gpt3}.
O FasiTech utiliza um prompt de \textit{few-shot} implícito com
especificação de:

\begin{itemize}
  \item \textbf{Papel e tarefa} — ``Você é um especialista em extrair
        informações de certificados acadêmicos''.
  \item \textbf{Formatos esperados} — lista de padrões comuns
        (``40h'', ``40 horas'', ``CH: 40h'', etc.).
  \item \textbf{Formato de saída} — estrutura canônica por página e
        linha \texttt{TOTAL GERAL: X horas}.
  \item \textbf{Tratamento de casos-limite} — carga horária em dias
        (8h/dia), múltiplas cargas na página, ausência de carga.
\end{itemize}

\subsection{Extração do Total por Expressão Regular}

Após a resposta do modelo, o sistema executa parsing determinístico
para extrair o total numérico, mitigando alucinações:

\begin{lstlisting}[language=Python,caption={Extração de total com regex.}]
padroes = [
    r'TOTAL GERAL:\s*(\d+(?:,\d+)?)\s*horas?',
    r'Total Geral:\s*(\d+(?:,\d+)?)\s*horas?',
    r'TOTAL:\s*(\d+(?:,\d+)?)\s*horas?',
]
for padrao in padroes:
    match = re.search(padrao, texto, re.IGNORECASE)
    if match:
        horas = match.group(1).replace(',', '.')
        return f"TOTAL GERAL: {horas} horas"
\end{lstlisting}

\section{Análise de Confiabilidade e Limitações}

\subsection{Fontes de Erro}

\begin{table}[H]
\centering
\caption{Fontes de erro e mitigações no módulo ACC.}
\label{tab:erros}
\begin{tabularx}{\textwidth}{lXX}
\toprule
\textbf{Fonte} & \textbf{Descrição} & \textbf{Mitigação} \\
\midrule
Baixa resolução & PDFs escaneados em $<150$ DPI & Reprocessamento a 300 DPI \\
Tabelas complexas & Layout de certificado não padrão & Prompt com múltiplos formatos \\
Alucinação do modelo & LLM inventa carga não presente & Regex pós-geração + fallback manual \\
Idioma estrangeiro & Certificados em inglês/espanhol & Instrução multilíngue no prompt \\
\bottomrule
\end{tabularx}
\end{table}

\subsection{Escalabilidade}

O modelo é chamado uma vez por envio (PDF completo), tornando o custo
fixo por aluno independente do número de páginas dentro do limite de
contexto do Gemini 2.5 Flash Lite ($\approx 1{,}5\times10^6$ tokens).
Para envios com muitos certificados, o sistema pode ser estendido com
paginação em lotes de 10--15 páginas por chamada.

% ─────────────────────────────────────────────────────────────────────────────
\chapter{Infraestrutura e Integração}
% ─────────────────────────────────────────────────────────────────────────────

\section{Banco de Dados Vetorial: LanceDB}

LanceDB \autocite{lancedb2023} é um banco de dados vetorial embutido
(\textit{embedded}) baseado no formato Apache Lance, projetado para
alto desempenho em hardware de commodity.
No FasiTech, ele persiste:

\begin{itemize}
  \item Tabela \texttt{recipes} — chunks dos documentos com vetores 768-d.
  \item Tabela \texttt{semantic\_cache} — entradas de cache semântico com
        vetores, metadados e timestamps.
\end{itemize}

\section{Modelos de Linguagem e Embedder}

O FasiTech utiliza o framework \textbf{Agno} (ex-Phidata) para abstrair
a camada de LLM. A seleção do modelo de geração é feita em tempo de
inicialização do \texttt{ChatbotService}:

\begin{enumerate}
  \item Se \texttt{GOOGLE\_API\_KEY} estiver configurada: usa
        \textbf{Gemini} via API Google.
  \item Caso contrário: usa \textbf{Ollama} \autocite{ollama2023} local como fallback
        (modelo \texttt{OLLAMA\_LLM\_MODEL}).
\end{enumerate}

\begin{table}[H]
\centering
\caption{Configuração dos modelos no FasiTech (\texttt{LLMConfig.py}).}
\begin{tabularx}{\textwidth}{lXX}
\toprule
\textbf{Variável} & \textbf{Uso} & \textbf{Padrão} \\
\midrule
\texttt{GEMINI\_MODEL} & Geração de respostas — Diretor Virtual & \texttt{gemini-2.5-flash} \\
\texttt{GEMINI\_MODEL\_VISION} & Extração de certificados — ACC & \texttt{gemini-2.5-flash-lite} \\
\texttt{OLLAMA\_LLM\_MODEL} & Fallback local (sem API key) & \texttt{qwen2.5:3b} \\
\texttt{nomic-embed-text} & Embeddings 768-d (sempre Ollama) & \texttt{OllamaEmbedder} \\
\bottomrule
\end{tabularx}
\end{table}

\noindent\textbf{Nota:} O embedder (\texttt{nomic-embed-text}) é sempre local via Ollama,
independentemente da disponibilidade da API Google. O modelo de geração
(Gemini vs.\ Ollama) é o único componente que muda conforme o ambiente.

\section{Padrão de Projeto: Singleton com Inicialização Lazy}

O \texttt{ChatbotService} implementa o padrão Singleton via \texttt{\_\_new\_\_},
garantindo que a inicialização custosa (conexão com Ollama, criação do
índice LanceDB, carregamento dos documentos) ocorra apenas uma vez por
processo. O diretório de dados é \texttt{\textasciitilde/.cache/fasitech/rag/}
(persistência por usuário do SO, compatível com ambientes Docker).

Sessões de usuário são isoladas por \texttt{session\_id}; cada sessão
cria um \texttt{Agent} Agno sob demanda com histórico de conversa
independente armazenado em SQLite (\texttt{.db} no mesmo diretório de cache).
Quando \texttt{persist\_history=False}, o histórico fica apenas em RAM.

% ─────────────────────────────────────────────────────────────────────────────
\chapter{Trabalhos Futuros}
% ─────────────────────────────────────────────────────────────────────────────

\begin{enumerate}
  \item \textbf{Multi-Query Retrieval} — gerar 3 variações da pergunta via
        LLM antes da busca e fundir resultados com RRF
        \autocite{es2023ragas,langchain2023multiquery}.
  \item \textbf{HyDE} (\textit{Hypothetical Document Embeddings}) — embeddar
        uma resposta hipotética gerada pelo LLM em vez da query original
        para maior precisão semântica \autocite{gao2022hyde}.
  \item \textbf{Parent-Child Chunking} — indexar chunks pequenos para
        busca, retornar chunk pai para geração \autocite{llamaindex2023}.
  \item \textbf{Cross-Encoder Reranking} — aplicar um reranker
        (\texttt{BAAI/bge-reranker-base}) sobre os chunks recuperados
        antes de montar o prompt \autocite{khattab2020colbert}.
  \item \textbf{Interface de Feedback} — expor botões de avaliação
        (1--5 estrelas) no ChatWidget para alimentar o cache semântico
        com dados reais de usuários.
  \item \textbf{Validação Cruzada ACC} — comparar resultado do Gemini
        com extração por OCR clássico (Tesseract) para medir
        concordância e detectar certificados problemáticos.
\end{enumerate}

% ─────────────────────────────────────────────────────────────────────────────
\printbibliography[title={Referências}]
% ─────────────────────────────────────────────────────────────────────────────

\end{document}
